% =====================================================================
%  APPENDIX B : External tools, explained from scratch
% =====================================================================
\chapter{External Tools, Explained from Scratch}
\label{app:tools}

\Daedalus{} is glue around a handful of powerful third-party libraries.  If you
came to this code from the physics side, the unfamiliar part is rarely the
field theory --- it is the \emph{tooling}: what \emph{is} SageMath, why does
everything run under \code{sage -python}, what does \code{nauty} compute, why do
the simulators use \code{numba} but the symbolic pipeline does not?  This
appendix answers those questions one tool at a time.  For each library we say
\textbf{what it is}, \textbf{why \Daedalus{} needs it}, and \textbf{where in the
code it appears} (with the rough usage count found by grepping the source, so
you can gauge how central it is).

\begin{note}
You do not install any of these by hand for the symbolic pipeline.  They all
come bundled inside the SageMath distribution (see the project's
\code{MSRJD\_diagrams} conda environment).  The simulators add \code{numba},
which is a plain \code{pip}/\code{conda} package.
\end{note}

% ---------------------------------------------------------------------
\section{SageMath --- the symbolic backbone}
% ---------------------------------------------------------------------

\textbf{What it is.}  SageMath (``Sage'') is a large open-source
\term{computer algebra system} (CAS) written in Python.  A CAS manipulates
mathematics \emph{symbolically} --- it knows that $\partial_t e^{-\mu t}=-\mu
e^{-\mu t}$ as an identity, not as a floating-point approximation on a grid.
Sage is really a Python front-end that stitches together dozens of specialised
mathematical engines (Pari for number theory, Singular for polynomials, GAP for
group theory, Maxima for calculus, \code{nauty} for graphs, \code{mpmath} for
arbitrary precision, and many more) behind one uniform Python API.  Because the
front-end \emph{is} Python, a Sage program is a Python program with extra
built-in types.

\textbf{Why \Daedalus{} needs it.}  Every object in the front half of the
pipeline is a symbol, not a number: the action $S$, the inverse propagator
$K_{\mathrm{ft}}(\omega)$ and its matrix inverse $G_{\mathrm{ft}}$, the
interaction vertices, the loop kernels.  These must be differentiated, inverted,
Fourier-transformed, Taylor-expanded, and have their poles found \emph{exactly}.
That is exactly what a CAS is for.

\begin{defn}
The \term{Symbolic Ring}, written \code{SR} in the code, is the Sage ``parent''
that all symbolic expressions belong to --- the analogue of \code{float} for
numbers.  A Sage symbol like \code{var('mu')} is an element of \code{SR}, and any
formula built from it (\code{mu + I*omega}, \code{exp(-mu*t)}) is again an
\code{SR} element.
\end{defn}

\subsection{Running under \texttt{sage -python}}
Because the pipeline imports \code{sage.all}, it cannot run under a stock Python
interpreter --- it must run under the Python that ships \emph{inside} Sage.  That
is why throughout this project you launch scripts as \code{sage -python
myscript.py} rather than \code{python myscript.py}, and why the notebooks use a
Sage Jupyter kernel.

\subsection{Parsing model text: \texttt{preparse} and \texttt{sage\_eval}}
When you type an action string such as
\code{sum(xt[i]*((Dt+mu)*x[i] + eps*x[i]\textasciicircum 3) - D*xt[i]\textasciicircum 2 for i in pop)},
two Sage facilities turn that text into a live symbolic object
(see Chapter~\ref{ch:theory-spec}):
\begin{description}[leftmargin=1.2em]
  \item[\code{preparse}] (used $\sim$7 places) rewrites Sage's mathematician-friendly
        shorthand into legal Python --- most importantly it turns the caret
        \code{\textasciicircum} (power) into \code{**}, and wraps integer
        literals so that \code{1/3} is the exact rational $\tfrac13$ rather than
        floating-point $0.333\ldots$
  \item[\code{sage\_eval}] (used $\sim$10 places) evaluates the preparsed string in a
        controlled namespace where the field names, parameters, and operator
        symbols (\code{Dt}, \code{Lap}) are already defined, returning an
        \code{SR} expression.
\end{description}

\subsection{Calculus, inversion, and poles}
\begin{itemize}
  \item \textbf{Matrix inversion} gives $G_{\mathrm{ft}}=K_{\mathrm{ft}}^{-1}$
        symbolically (Chapter~\ref{ch:propagator}).
  \item \code{fourier\_transform} (used $\sim$8 places) moves kernels between the
        time and frequency domains.
  \item \code{.roots()} (used $\sim$6 places) solves the characteristic
        polynomial of the denominator to locate the \emph{retarded poles}
        (those with $\operatorname{Im}\omega>0$) that drive the time-domain
        propagator.
  \item \code{assume(...)} declares domain facts (e.g.\ \code{mu>0}) so the CAS
        can simplify $|{\cdot}|$, choose branch cuts, and evaluate integrals
        unambiguously.
\end{itemize}

\subsection{Graphs inside Sage (and the door to \texttt{nauty})}
Sage also has first-class \term{graph} objects.  \Daedalus{} builds diagram
topologies as \code{Graph} / \code{DiGraph} and then calls
\code{canonical\_label}, \code{is\_isomorphic}, and \code{automorphism\_group} on
them.  Those three methods are the bridge to \code{nauty} (next section): they
are how the enumerator decides that two diagrams are ``the same'' and how it
counts a diagram's symmetries.

\subsection{\texttt{.sobj} --- Sage's saved-object format}
A \code{.sobj} file is Sage's version of a Python \code{pickle}: a binary dump of
a live Sage object.  The enumeration cache writes the (expensive-to-compute)
list of unique typed diagrams to \file{saved\_theories/<name>/...\_.sobj} so the
next run loads them in milliseconds instead of re-enumerating
(Chapter~\ref{ch:caching}).

% ---------------------------------------------------------------------
\section{nauty --- canonical graphs and automorphisms}
% ---------------------------------------------------------------------
\textbf{What it is.}  \code{nauty} (a playful contraction of ``\emph{No
AUTomorphisms, Yes?}''), by Brendan McKay, is the standard high-performance tool
for two closely-related graph problems:
\begin{description}[leftmargin=1.2em]
  \item[Canonical labelling] relabel a graph's vertices into a unique normal
        form so that two graphs are \term{isomorphic} (the same up to renaming
        vertices) if and only if their canonical forms are byte-identical.  This
        turns the hard ``are these two diagrams the same?'' question into a
        cheap string comparison.
  \item[Automorphism group] compute $\Aut(\Gamma)$, the group of vertex
        relabellings that map the graph to itself --- i.e.\ its symmetries ---
        and in particular its size $|\Aut(\Gamma)|$.
\end{description}

\textbf{Why \Daedalus{} needs it.}  Two things in diagrammatics reduce to graph
symmetry:
\begin{enumerate}
  \item \textbf{De-duplication.}  The raw enumerator produces many topologies
        that are really the same diagram drawn differently; canonical labelling
        collapses them to one representative (Chapter~\ref{ch:enumeration}).
  \item \textbf{The symmetry factor $\Sym(\Gamma)$.}  Each Feynman diagram is
        divided by the order of its automorphism group (with external legs held
        fixed).  \Daedalus{} computes this exactly from
        $|\Aut(\Gamma)|$ rather than guessing a combinatorial fudge factor
        (Chapter~\ref{ch:type-assignment}).
\end{enumerate}

\begin{gotcha}
You never call \code{nauty} directly in this codebase --- it is reached
\emph{transitively} through Sage's graph methods.  If you grep for ``nauty'' you
will find almost nothing; look for \code{automorphism\_group} and
\code{canonical\_label} instead.
\end{gotcha}

% ---------------------------------------------------------------------
\section{NumPy --- arrays and fast numerics}
% ---------------------------------------------------------------------
\textbf{What it is.}  NumPy is the foundational numerical-array library for
Python: the \code{ndarray} (an $n$-dimensional, homogeneously-typed array) plus
vectorised operations on it.  \textbf{Where it appears.}  Everywhere the pipeline
turns symbols into numbers: the lag grid $\tau$, the correlator arrays
$C(\tau)$, the residue matrices, the spatial momentum grids, and the simulator
output bins are all \code{ndarray}s.  When this manual says ``a callable returns
an array of shape \code{(k-1, n\_tau)}'', that array is a NumPy array.

% ---------------------------------------------------------------------
\section{SciPy --- integration, linear algebra, special functions}
% ---------------------------------------------------------------------
\textbf{What it is.}  SciPy is the scientific-computing layer built on top of
NumPy.  \Daedalus{} uses four of its sub-packages:
\begin{description}[leftmargin=1.2em]
  \item[\code{scipy.integrate}] adaptive numerical integration.  \code{quad} does
        a single one-dimensional integral; \code{nquad} does a nested
        multi-dimensional one.  These are the \emph{fallback} evaluators in the
        temporal Phase~J engine (used when the analytic mode-sum does not apply)
        and appear in the spatial chamber integrals
        (Chapters~\ref{ch:phasej-temporal},~\ref{ch:spatial-core}).
  \item[\code{scipy.linalg}] dense linear algebra.  \code{expm} is the
        matrix exponential $e^{At}$ (time-domain propagation);
        \code{solve\_continuous\_lyapunov} solves the Lyapunov equation
        $AX+XA^{\!\top}=-Q$ for a stationary covariance; eigenvalue routines back
        the linear-stability test (Chapter~\ref{ch:mean-field}).
  \item[\code{scipy.special}] special functions.  \code{kv} is the
        \term{modified Bessel function of the second kind} $K_\nu(z)$, which
        emerges from the radial Schwinger-parameter integral in the spatial
        backend; \code{gamma} is Euler's $\Gamma$ function
        (Chapter~\ref{ch:spatial-heatkernel}).
  \item[\code{scipy.optimize}] root finding.  \code{root} / \code{fsolve} are
        used as a fallback for the mean-field self-consistency solve
        (Chapter~\ref{ch:mean-field}).
\end{description}

% ---------------------------------------------------------------------
\section{sympy and \texttt{lambdify}}
% ---------------------------------------------------------------------
\textbf{What it is.}  \code{sympy} is a second, \emph{pure-Python} computer
algebra system (independent of Sage).  Its \code{lambdify} function takes a
symbolic expression and returns a fast ordinary Python function that evaluates it
on NumPy arrays --- it ``compiles'' a formula into a numeric callable.
\textbf{Where it appears.}  Narrowly: a small number of modules convert a
symbolic spatial form-factor into a \code{lambdify}-ed callable so it can be
sampled cheaply on a momentum grid.  The bulk of the symbolic work stays in
Sage; \code{sympy} is used only at this symbolic-to-numeric hand-off.

% ---------------------------------------------------------------------
\section{numba --- JIT compilation for the simulators}
% ---------------------------------------------------------------------
\textbf{What it is.}  \code{numba} is a just-in-time (JIT) compiler: decorate a
numeric Python function with \code{@numba.njit} and the first call compiles it to
machine code through LLVM, after which it runs at C-like speed.
\textbf{Where it appears.}  \emph{Only} in the independent simulators
(\file{models/}, $\sim$11 files), never in the symbolic pipeline.  The
Euler--Maruyama and Hawkes integrators are tight inner loops over millions of
time-steps --- exactly the workload JIT compilation is for
(Chapter~\ref{ch:simulators}).

\begin{gotcha}
\code{numba} can only compile functions whose arguments are plain machine types
(\code{float}, \code{int}, NumPy arrays).  The resolved parameters coming out of
the pipeline are \emph{Sage ring elements}, which \code{numba} cannot type ---
so the simulator-driving notebook cells explicitly cast every scalar with
\code{float(...)} before calling an \code{@njit} simulator.  Forgetting the cast
gives an opaque numba typing error.
\end{gotcha}

% ---------------------------------------------------------------------
\section{matplotlib --- plotting}
% ---------------------------------------------------------------------
\textbf{What it is.}  The standard Python 2-D plotting library.
\textbf{Where it appears.}  \code{dd.plot\_cumulant} builds the theory-vs-lag
curves (and the $k\!\ge\!3$ slice/heat-map panels and the simulator overlays)
with matplotlib \code{Figure}/\code{Axes} objects
(Chapter~\ref{ch:engine-api}).

% ---------------------------------------------------------------------
\section{networkx --- a lightweight graph helper}
% ---------------------------------------------------------------------
\textbf{What it is.}  \code{networkx} is a pure-Python graph library (nodes,
edges, traversals).  \textbf{Where it appears.}  In a single module, as a
convenience for representing/handling diagram structure; the heavy graph work
(canonicalisation, automorphisms) is done by Sage/\code{nauty}, not
\code{networkx}.

% ---------------------------------------------------------------------
\section{multiprocessing --- parallelism, and the macOS fork hazard}
% ---------------------------------------------------------------------
\textbf{What it is.}  Python's standard library for running work in separate
\emph{processes} (sidestepping the Global Interpreter Lock).  A process pool can
\emph{fork} (clone the current process) or \emph{spawn} (start a fresh
interpreter).  \textbf{Where it appears.}  $\sim$7 modules optionally parallelise
diagram enumeration and the Phase~J $\tau$-batch with a fork-based pool.

\begin{gotcha}
On macOS, \code{fork}-ing a process that has already initialised Cocoa, BLAS, or
a Matplotlib backend --- which a Jupyter kernel has --- can hard-crash the kernel
\emph{and} the OS.  \Daedalus{} therefore guards every fork behind a
\term{fork-safety} check that degrades to serial execution inside a macOS Jupyter
kernel, while keeping the fork (the fast path) in terminals, scripts, and on
Linux.  This is covered in detail in Chapter~\ref{ch:type-assignment}.
\end{gotcha}

% ---------------------------------------------------------------------
\section{mpmath --- arbitrary-precision arithmetic}
% ---------------------------------------------------------------------
\textbf{What it is.}  \code{mpmath} is a pure-Python library for real and complex
floating-point arithmetic at \emph{arbitrary} precision (hundreds of digits if
asked).  Sage uses it internally, and a few precision-sensitive spots in the
pipeline ($\sim$3 modules) lean on it where ordinary double precision would lose
significant digits --- for example in mode-sums with closely-spaced poles
(see the $m\!\ge\!3$ precision note in Chapter~\ref{ch:phasej-temporal}).

% ---------------------------------------------------------------------
\section{The Python standard library}
% ---------------------------------------------------------------------
A few standard modules do quiet but essential work:
\begin{description}[leftmargin=1.2em]
  \item[\code{ast}] safe parsing of Python literals; \code{ast.literal\_eval}
        reads the \code{METADATA}/\code{DEFAULT\_FUNDAMENTAL} dictionaries out of
        a \file{.theory.py} file without executing arbitrary code
        (Chapter~\ref{ch:theory-files}).
  \item[\code{importlib}] dynamic import; loads a \file{.theory.py} file by path
        at run time so new theories need no registration.
  \item[\code{re}] regular expressions; light text munging of action/equation
        strings.
  \item[\code{dataclasses}] the \code{@dataclass} decorator that defines the
        \code{Config} object (Chapter~\ref{ch:engine-api}).
  \item[\code{itertools}] combinatorial iterators; \code{product} generates leg
        assignments and Wick pairings.
\end{description}

\begin{note}
\textbf{Rule of thumb for ``what library does this?''}  Anything
\emph{symbolic} (expressions, inversion, Fourier transforms, graph symmetry) is
\textbf{Sage} (with \code{nauty} under the graph hood).  Anything
\emph{numeric in the pipeline} (integrals, eigenvalues, Bessel functions) is
\textbf{SciPy/NumPy}.  Anything in a \emph{simulator} is \textbf{numba+NumPy}.
The \emph{UI} is \textbf{ipywidgets}.  Everything else is the Python standard
library.
\end{note}
